
\minitoc

% ========================================================================================================
% Introduction 

Nous définissons dans ce chapitre un nouveau type de processus aléatoires plus théorique 
que les chaînes de Markov. Ce sera un outil mathématique du fait de ses bonnes 
propriétés de convergence. Nous appelerons ces processus des Martingales dont l'espérance 
conditionnelle par rapport au passé restera constante. 

\medskip 

Dans tout ce chapitre, nous considérerons des processus aléatoires à valeurs dans $\R$. 

% ========================================================================================================
% Généralités

\section{Généralités}

\subsection{Définition et exemples}

\begin{definition}[Martingale]
    Soit $( \Omega, \mathcal{A}, \mathcal{F}, \myP)$ un espace de probabilité filtré 
    et $(X_n)_{n \in \N}$ un processus aléatoire intégrable. On dira que $(X_n)_{n \in \N}$ est 
    \begin{enumerate}[label=\roman*)]
        \item une \emph{martingale} par rapport à $ (\mathcal{F}_n)_{n \in \N}$ si : 
            \begin{equation}\label{definition:martingale}
                \forall n \geqslant 1, \quad \mathbb{E}(X_n \mid \mathcal{F}_{n-1}) = X_{n-1}
            \end{equation}
        \item une \emph{sur-martingale} par rapport à $ (\mathcal{F}_n)_{n \in \N}$ si : 
            \begin{equation}
                \mathbb{E}(X_n \mid \mathcal{F}_{n-1}) \leqslant X_{n-1}
            \end{equation}
        \item une \emph{sous-martingale} par rapport à $ (\mathcal{F}_n)_{n \in \N}$ si : 
            \begin{equation}
                \mathbb{E}(X_n \mid \mathcal{F}_{n-1}) \geqslant X_{n-1}
            \end{equation}
    \end{enumerate}
\end{definition}

\begin{example}[Marche aléatoire simple]
    Soit $( \varepsilon_n)_{n \in \N}$ une suite indépendantes de variables aléatoires 
    à valeurs dans $(E, \mathcal{E})$ telles que $ \mathbb{E}( \varepsilon_n) = 0$ pour tout $n$. 
    On définit une marche aléatoire sur $E$ comme la suite $(S_n)_{n \in \N}$ telle que : 
        $$ S_0 = 0 \quad \forall n \geqslant 1, \; S_n = \sum_{k=1}^{n} \varepsilon_k $$ 
    $(S_n)_{n \in \N}$ est une martingale pour la filtration $ \mathcal{F}_n = \sigma( \varepsilon_1, \dots, \varepsilon_n)$. 
    En effet, on a : 
        \begin{align*}
            \mathbb{E}(S_{n+1} \mid \mathcal{F}_n) &= \mathbb{E}(S_n + \mathcal{E}_{n+1} \mid \mathcal{F}_n) \\ 
            &= \mathbb{E}(S_n \mid \mathcal{F}_n) + \mathbb{E}( \varepsilon_{n+1} \mid \mathcal{F}_n) 
        \end{align*}
    Or $ \varepsilon_{n+1} \perp \varepsilon_0, \dots, \varepsilon_n$ donc $ \mathbb{E}( \varepsilon_{n+1} \mid \mathcal{F}_n) = 0 $ d'où 
        $$ \mathbb{E}( S_{n+1} \mid \mathcal{F}_n) = \mathbb{E}(S_n) = S_n $$ 
    c'est donc bien une martingale par définition. On peut remarquer que si $ \mathbb{E}(S_n) \geqslant 0$, alors 
    $(S_n)_{n \in \N}$ est une sous-martingale. 
\end{example}


\begin{remark}
    Si $ \mathbb{E}( \varepsilon_n) \geqslant 0, \forall n \in \N$ alors : 
        $$ \mathbb{E}( S_{n+1} \mid \mathcal{F}_n) = S_n + \mathbb{E}( \varepsilon_{n+1}) \geqslant S_n $$ 
    donc $(S_n)_{n \in \N}$ serait une sous-martingale. 
\end{remark}

\begin{example}[Multiplicatif]
    Soit $( \varepsilon_i)_{i \in \N}$ une suite indépendante de variables aléatoires telles que $ \mathbb{E}( \varepsilon_i) = 1$ pour 
    tout $ i \in \N$. Posons : 
        $$ M_0 = 1 \quad \text{et} \quad \forall n \in \N, \; M_n = \prod_{i = 1}^{n} \varepsilon_i $$ 
    et $ \mathcal{F}_n = \sigma ( \varepsilon_1, \dots, \varepsilon_n)$. Montrons que $(M_n)_{n \in \N}$ est une 
    martingale. Pour tout $ n \in \N$, on a : 
        \begin{align*}
            \mathbb{E}(M_{n+1} \mid \mathcal{F}_n) &= \mathbb{E} \left( \varepsilon_{n+1} \times M_n \mid \mathcal{F}_n\right) \\ 
            &= M_n \times \mathbb{E}( \varepsilon_{n+1} \mid \mathcal{F}_n) \\
            &= M_n \times \mathbb{E}( \varepsilon_{n+1}) \\ 
            &= M_n 
        \end{align*}
    Par définition, $(M_n)_{n \in \N}$ est donc une martingale. 
\end{example}

\begin{example}[Lien avec les chaînes de Markov]
    On peut généraliser les chaînes de Markov en tant que Martingales grâce à l'utilisation de fonctions harmoniques. 
    Soit $(X_n)_{n \in \N}$ sur un espace d'états discret $E$ et de matrice de transition $P$. Une fonction 
    harmonique $h : E \longrightarrow \R$ vérifie : 
        $$ \forall x \in E, \quad h(x) = \sum_{y \in E} P(x,y)h(y) $$ 
    c'est donc un \emph{vecteur propre à droite} pour $P$. Supposons de plus que $ \forall n \in \N, h(X_n) \in L^1(\R)$. 
    Alors $(h(X_n))_{n \in \N}$ sera une martingale pour la filtration $( \mathcal{F}_n)_{n \in \N}$ associée à $(X_n)_{n \in \N}$. 
    En effet, pour tout $ n \in \N$, on a : 
        \begin{align*}
            \mathbb{E}(h(X_{n+1}) \mid \mathcal{F}_n) &= \mathbb{E}(h(X_{n+1}) \mid X_n) \\ 
            &= \sum_{y \in E} \mathbb{E}(h(X_{n+1}) \mathbbm{1}_{X_{n+1} = y} \mid X_n ) \\ 
            &= \sum_{y \in E} \myP(X_{n+1} = y \mid X_n) \\ 
            &= \sum_{y \in E} P(X_n,y) \\ 
            &= h(X_n)
        \end{align*}
\end{example}

\begin{proposition}[Généralisation de la définition]
    Soit $(M_n)_{n \geqslant 0}$ une martingale pour une filtration $ \mathcal{F}_n$ alors : 
        \begin{equation}\label{proposition:generalisation-definition}
            \forall n \in \N, \forall k \geqslant 1, \quad \boxed{ \mathbb{E}(M_{n +k} \mid \mathcal{F}_n) = M_n}
        \end{equation}
\end{proposition}

\begin{proof}
    Par récurrence sur $k$, on a :
    \begin{itemize}
        \item \textbf{Initialisation} : pour $k = 1$ c'est la définition. 
        \item \textbf{Hérédité} : soit $k \geqslant 1$ tel que : 
            $$ \mathbb{E}(M_{n+k} \mid \mathcal{F}_n) = M_n$$ 
        alors : 
            \begin{align*}
                \mathbb{E}(M_{n + k + 1} \mid \mathcal{F}_n) &= \mathbb{E}( \mathbb{E}(M_{n+k+1} \mid \mathcal{F}_n) \mid \mathcal{F}_n) \\ 
                &= \mathbb{E}(M_{n+1} \mid \mathcal{F}_n) = M_n 
            \end{align*}
    \end{itemize}
\end{proof}

\begin{remark}
    La proposition précédente nous donne une relation de récurrence pour les espérances : 
        $$ \forall n \in \N, \forall k \geqslant 1, \quad \mathbb{E}( \mathbb{E}(M_{n+k} \mid \mathcal{F}_n)) = \mathbb{E}(M_n) $$ 
    or $ \mathbb{E}( \mathbb{E}(M_{n+k} \mid \mathcal{F}_n)) = \mathbb{E}(M_{n+k})$ d'où : 
        $$ \mathbb{E}(M_{n+k}) = \mathbb{E}(M_n) $$ 
        $$ \iff \boxed{ \mathbb{E}(M_n) = \mathbb{E}(M_0)} $$ 
    Nous essaierons par la suite de généraliser cette proposition pour des temps aléatoires. 
\end{remark}

\subsection{Lien Martingale et sous-martingale}

Pour rappel, une sous-martingale est un processus $(Y_n)$ adapté à la filtration $ \mathcal{F}_n$ tel que : 
    $$ \mathbb{E}(Y_{n+1} \mid \mathcal{F}_n) \geqslant Y_n $$ 
De plus, pour une fonction convexe $g : \R \longrightarrow \R$, pour $X \in L^1( \mathcal{A}, \myP)$
et $ \mathcal{F}$ une sous-tribu de $ \mathcal{A}$, l'inégalité de Jensen nous donne 
    $$ \mathbb{E}(g(X) \mid \mathcal{F}) \geqslant g( \mathbb{E}(X \mid \mathcal{F})) $$ 
Nous pouvons donc faire le lien entre les martingales et les sous-martingales. 

\begin{corollary}[Inégalité de Jensen]\label{corollary:inegalite-jensen}
    Soit $(M_n)_{n \in \N}$ une martingale pour une filtration $ (\mathcal{F}_n)_{n \geqslant 0}$. 
    Soit $g : \R \longrightarrow \R$ \emph{convexe} telle que $g(M_n) \in L^1$ pour tout $n \in \N$.
    Alors $(g(M_n))_{n \in \N}$ est une sous-martingale. 
\end{corollary}

\begin{proof}
    On a facilement $ \mathbb{E}(g(M_{n+1}) \mid \mathcal{F}_n) \geqslant g( \mathbb{E}(M_{n+1} \mid \mathcal{F}_n)) = g(M_n)$. 
\end{proof}

\begin{example}[Marche aléatoire]
    Soit $( \varepsilon_i)_{i \geqslant 0}$ une suite idd de variables aléatoires telles que : 
        $$ \forall i \in \N, \quad \mathbb{E}( \varepsilon_i) = 0 \; \text{ et } \; \varepsilon_i^2 \in L^1 $$ 
    Posons $ S_n = \sum_{i=1}^{n} \varepsilon_i$ et $ \mathcal{F}_n = \sigma( \varepsilon_1, \dots \varepsilon_n)$. 
    $(S_n)_{n \geqslant 0}$ est donc une martingale pour $ \mathcal{F}_n$ (voir exemple précédent). D'après le 
    corollaire de l'inégalité de Jensen, $(S_n^2)_{n \geqslant 0}$ est une sous-martingale. On a donc : 
        $$ \mathbb{E}(S_{n+1}^2 \mid \mathcal{F}_n) \geqslant \mathbb{E}(S_n^2) $$ 
        $$ \iff \mathbb{E}(S_{n+1}^2) \geqslant \mathbb{E}(S_n^2) $$ 
    On a donc une suite croissante d'espérances et : 
        $$ \mathbb{E}(S_n^2) = \mathbb{E} \left( \left(\sum_{i=1}^{n} \varepsilon_i\right)^2 \right) = n \mathbb{E}( \varepsilon_1 ^2) $$ 
    Essayons donc de construire une martingale en posant : 
        $$ \forall n \in \N, \quad M_n = S_n^2 - \mathbb{E}(S_n^2) = S_n^2 - n \mathbb{E}( \varepsilon_1^2) $$ 
    On a donc, par définition : 
        $$ \mathbb{E}(M_{n+1} \mid \mathcal{F}_n) = \mathbb{E}(S_{n+1}^2 \mid \mathcal{F}_n) - (n+1)\mathbb{E}( \mathbb{E} (\varepsilon_1^2) \mid \mathcal{F}_n) $$ 
    or  
        \begin{align*}
            \mathbb{E}(S_{n+1}^2 \mid \mathcal{F}_n) &= \mathbb{E}\big[ ( S_n + \varepsilon_{n+1})^2 \mid \mathcal{F}_n \big] \\ 
            &= \mathbb{E}(S_n^2 \mid \mathcal{F}_n) + 2\mathbb{E}(S_n \varepsilon_{n+1} \mid \mathcal{F}_n) + \mathbb{E}( \varepsilon_{n+1}^2 \mid \mathcal{F}_n) \\ 
            &= S_n^2 + 2 S_n \mathbb{E}( \varepsilon_{n+1}) + \mathbb{E}( \varepsilon_{n+1}^2) \\ 
            &= S_n^2 + \mathbb{E}( \varepsilon_1^2) 
        \end{align*}
    d'où : 
        \begin{align*}
            \mathbb{E}(M_{n+1} \mid \mathcal{F}_n) = S_n^2 - n \mathbb{E}( \varepsilon_1^2) = M_n 
        \end{align*}
    On peut donc décomposer notre sous-martingale $S_n^2$ en une somme de martingale et ce que l'on appelera 
    un processus croissant en $n$ : 
        $$ S_n^2 = M_n + n \mathbb{E}( \varepsilon_1^2) $$ 
\end{example}

\begin{proposition}[Décomposition de Doob]
    Soit $(X_n)_{n \geqslant 0}$ une sous-martingale par rapport à une filtration $ \mathcal{F}_n$. 
    Alors il existe un unique couple de processus $(M_n)_{n \geqslant 0}$ et $(Z_n)_{n \geqslant 0}$ tels que: 
    \begin{enumerate}[label=\roman*)]
        \item $(M_n)_{n \geqslant 0}$ est une martingale pour $ \mathcal{F}_n$ 
        \item $Z_0 = 0$ et $Z_{n+1} \geqslant Z_n$ et $(Z_n)_{n \geqslant 0}$ est prévisible pour $ (\mathcal{F}_n)_{n \geqslant 0}$. 
    \end{enumerate}
    et : 
        \begin{equation}\label{eq:doob-inequality}
            \forall n \in \N, \quad \boxed{X_n = M_n + Z_n}
        \end{equation}
\end{proposition}

\begin{proof}
    Soit $(X_n)_{n \in \N}$ une sous-martingale. par rapport à $( \mathcal{F}_n)_{n \geqslant 0}$. On définit les 
    écarts : 
        $$ \forall n \in \N, \quad \Delta_n = X_{n+1} - X_n $$ 
    Posons $ M_0 = X_0$ et $Z_0 = 0$. On construit récursivement $(M_n)_{n \geqslant 0}$ et $(Z_n)_{n \geqslant 0}$ tels 
    que :
        $$ \forall n \in \N, \quad Z_n = \sum_{i=1}^{n} \mathbb{E}( \Delta_i \mid \mathcal{F}_{i-1}), \quad M_n = X_n - Z_n $$
    Montrons que des deux processus vérifient les hypothèses : 
    \begin{enumerate}[label=\roman*)]
        \item Par contruction, pour tout $ n \in \N$, on a $X_n = M_n + Z_n$. 
        \item Pour tout $i \in \llbracket 1, n \rrbracket$, $ \mathbb{E}( \Delta_i \mid \mathcal{F}_{i-1})$ est $ \mathcal{F}_{i-1}$ 
            mesurable et $Z_n$ est une somme de $i$ variables $ \mathcal{F}_{i-1}$ mesurables. Donc $(Z_n)_{n \geqslant 0}$ est 
            prévisible pour $( \mathcal{F}_n)_{n \geqslant 0}$. 
        \item Pour tout $n \in \N$, on a : 
            \begin{align*}
                Z_{n+1} - Z_n &= \mathbb{E}( \Delta_{n+1} \mid \mathcal{F}_n) \\ 
                &= \mathbb{E}(X_{n+1} \mid \mathcal{F}_n) - \mathbb{E}(X_n \mid \mathcal{F}_n) \\ 
                &= \mathbb{E}(X_{n+1} \mid \mathcal{F}_n) - X_n 
            \end{align*}
        or $X_n$ est une sous-martingale donc : 
            $$ \mathbb{E}(X{n+1} \mid \mathcal{F}_n) \Longrightarrow Z_{n+1} \geqslant Z_n $$ 
        Puisque $Z_0 = 0$, $(Z_n)_{n \in \N}$ est un processus croissant. 
    \end{enumerate}
    Montrons maintenant que $(M_n)_{n \in \N}$ est une martingale. 
    \begin{align*}
        \forall n \in \N, \mathbb{E}(M_{n+1} \mid \mathcal{F}_n) &= \mathbb{E}(X_{n+1} \mid \mathcal{F}_n) - \mathbb{E}(Z_{n+1} \mid \mathcal{F}_n) \\ 
        &= \mathbb{E}(X_{n+1} \mid \mathcal{F}_n) - Z_{n+1} \quad \\ 
        &= \mathbb{E}(X_{n+1} \mid \mathcal{F}_{n}) - Z_n - \mathbb{E}( \Delta_{n+1} \mid \mathcal{F}_n) \\ 
        &= \mathbb{E}(X_{n+1} \mid \mathcal{F}_{n}) - Z_n - (\mathbb{E}( \Delta_{n+1} \mid \mathcal{F}_n) - X_n) \\ 
        &= X_n - Z_n = M_n
    \end{align*}
    L'unicité se montre par récurrence sur $n$. Nous ne la ferons pas ici.
\end{proof}


% ========================================================================================================
% Stratégies et théorème d'arrêt

\section{Stratégies et théorème d'arrêt}

\begin{example}[Ruine du joueur]
    Considérons un joueur lors d'un jeu aléatoire constitué de $N$ parties. Lors de la partie 
    $n \leqslant N$, il mise une somme $ \Theta_n$ et il remporte la some $ \varepsilon_n \Theta_n$ où 
    $ \varepsilon_n = -1$ s'il perd et $ \varepsilon_n = 2$ s'il gagne. 

    \medskip

    Les variables $( \varepsilon_n)_{n \geqslant 0}$ dépendent du jeu, donc on peut construire une filtration 
    $ ( \mathcal{F}_n)_{n \geqslant 0}$ définie par $ \mathcal{F}_n = \sigma( \varepsilon_1, \dots \varepsilon_n)$. 
    On peut considérer que $( \varepsilon_n)_{n \in \N}$ sont iid. Par construction, $\Theta_n$ est $\mathcal{F}_{n-1}$-mesurable, donc prévisible.
    On définit alors le gain total du joueur après $n$ parties comme la variable : 
        $$ \forall n \leqslant N, \quad S_n = S_0 + \sum_{k=1}^{n} \Theta_k \varepsilon_k $$ 
\end{example}

Dans cette section, nous allons essayer de généraliser cette structure au travers des stratégies. 

\begin{proposition}
    Soit $(M_n)_{n \geqslant 0}$ une martingale pour la filtration $( \mathcal{F}_n)_{n \geqslant 0}$, 
    $ ( \Phi_n)_{n \geqslant 0}$ un processus prévisible et borné. le processus défini par :
        \begin{equation}\label{eq:not-optimal-strategy}
            X_0 = 0 \text{ et } \forall n \in \N, \quad X_n = \sum_{k=1}^{n} \Phi_k (M_k - M_{k-1})
        \end{equation}
    Alors $(X_n)_{n \in \N}$ est une martingale pour $( \mathcal{F}_n)_{n \geqslant 0}$. 
\end{proposition}

Cette proposition est fondamentale dans l'interprétation des martingales en tant que stratégies de jeu. 
En effet, elle nous montre que pour n'importe quel jeu initial d'espérance nulle $(M_n)_{n \in \N}$ (gain moyen), il n'existe 
aucune stratégie $( \Phi_n)_{n \in \N}$ permettant de modifier le gain moyen. 
Cette notion est traduite par le fait que $(X_n)_{n \in \N}$ reste une martingale. 

\begin{remark}
    Les martingales nous permettent de modéliser la plupart des stratégies. Par exemple, si l'on veut 
    modéliser une stratégie du type "le joueur mise jusqu'à avoir accumulé une certaine somme $ \lambda$", 
    alors on peut définir $T$ le temps d'arrêt tel que : 
        $$ T = \inf \{n \in \N \mid S_n \geqslant \lambda\} $$ 
    On posera donc : 
        $$ \forall n \in \N, \forall k \in \llbracket 1, n \rrbracket, \quad \Phi_k = \mathbbm{1}_{k \geqslant \lambda} \quad \text{et} 
        \quad X_n = \sum_{k=1}^{n} \Phi_k (M_k - M_{k-1}) $$  
\end{remark}

On souhaiterai aussi généraliser la notion de chaîne de Markov $(X_n)_{n \in \N}$ arrêtée par un temps d'arrêt $T$ 
définie par : 
    $$ 
        \forall n \in \N, \quad X_n = 
            \begin{cases}
                X_n \quad \text{si } n \leqslant T \\ 
                X_T \quad \text{sinon}  
            \end{cases}
    $$

\begin{proposition}[Martingale arrêtée]
    Ainsi, soit $(M_n)_{n \leqslant 0}$ une martingale pour la filtration $ \mathcal{F}_n$ $T$. 
    Soit $T$ un temps d'arrêt. On pose : 
        $$ \forall k \in \N, \quad \Phi_k = \mathbbm{1}_{\{T \geqslant n\}} $$ 
    L'évènnement $\{T \leqslant n\}$ est mesurable par rapport à $ \mathcal{F}_{n-1}$ car il se décompose 
    comme le complémentaire d'une intersection dénombrable d'évènnements $ \mathcal{F}_{n-1}$ mesurables : 
        $$ \{T \leqslant n\} = \overline{\{T \geqslant n\}} = \bigcap_{k = 1}^{n-1} \{T \not = k\} $$ 
    Le processus $ \Pi_k $ est donc prévisible pour la filtration $( \mathcal{F}_n)_{n \geqslant 0}$. 
    D'après la proposition \ref{eq:not-optimal-strategy}, le processus défini comme : 
        \begin{equation}\label{eq:martingale-arrete}
            \boxed{
                \forall n \in \N, \quad M_n = 
            \begin{cases}
                M_n \quad \text{si } n \leqslant T \\ 
                M_T \quad \text{ sinon }  
            \end{cases}
            }
        \end{equation}
\end{proposition}

\begin{theorem}[Arrêt (Doob)]\label{theorem:stop-doob-theorem}
    Soit $(M_n)_{n \geqslant 0}$ une martingale pour la filtration $( \mathcal{F}_n)_{n \in \N}$.
    Soit $T$ un temps d'arrêt pour $( \mathcal{F}_n)_{n \geqslant 0}$. 
    Si une des trois assertions suivante est vraie : 
    \begin{enumerate}[label=\roman*)]
        \item $T$ est borné p.s 
        \item $ T < \infty$ p.s et $(M_n)_{n \in \N}$ est borné. 
        \item Si $T \in L^1$ et : 
            \begin{equation}\label{eq:third-stop-condition}
                \sup_n \left| M_n - M_{n-1 }\right| = \sup_n \left| \Delta_{M_n}\right| \leqslant C \quad \text{p.s}  
            \end{equation}
    \end{enumerate}
    Alors $\boxed{ \mathbb{E}(M_T) = \mathbb{E}(M_0)}$.
\end{theorem}

Ce théorème nous permet donc de généraliser aux temps d'arrêt la proposition de généralisation de la 
définition des martingales \ref{proposition:generalisation-definition}.